\section{Limits of Gadget Classes}
\label{sec:gadget-frontiers}

\Cref{sec:lp-hardness} made a local parity signal robust by
replicating constraints, scaling rows, or increasing the height of the
propagated signal.  This section asks whether some other sparse LP gadget
can avoid all three costs.  For several broad, precisely stated classes the
answer is no.  The most general result is a convex-mixture theorem that
assumes no graph structure.  For signed-copy graphs and bounded-arity Walsh
rows, electrical resistance gives a geometric refinement.  Conservation of
the dual multipliers shows that complete central and Newton outputs force a
large row scale, inverse dual scale, multiplier size, or input incidence.

Throughout this section the hidden input is
\(\varsigma\in\{\pm1\}^N\), with \(N\geq2\), and
\(p(\varsigma)=\prod_{i=1}^N\varsigma_i\).  Input locality always refers to
the raw sparse coefficient access of \cref{def:raw-sparse-access}.
A residual written as \(\norm{Ax-b}_2\) is absolute; when a relative
threshold is named, it is the global RHS-relative residual of
\cref{def:global-residual}.  The results are structural limits on
output requirements.  They become query lower bounds only after a separate
reduction shows that the oracle can be simulated locally and that parity can
be decoded.

\subsection{Convex mixtures of neighboring witnesses}
\label{sec:frontier-mixture}

Consider a fixed-support family
\begin{equation}
 \mathcal F_{\varsigma}
 =\{x\in\R_{\geq0}^n:A_{\varsigma}x=b\},
 \label{eq:frontier-local-family}
\end{equation}
where \(b\) is public.  Every row has at most \(s\) nonzeros, every
coefficient has magnitude at most \(B\), and each input-dependent position
\((r,j)\) is a function of one designated bit
\(\varsigma_{\ell(r,j)}\).  Let \(M_{\rm dep}\) be the number of these
positions.  A bit may label many positions, including several in one row.
Fix a base input, let \(\varsigma^{(i)}\) flip bit \(i\), and suppose that
each such neighbor has an exact witness
\begin{equation}
 x^{(i)}\in\mathcal F_{\varsigma^{(i)}},
 \qquad \norm{x^{(i)}}_\infty\leq H
 \quad(1\leq i\leq N).
 \label{eq:frontier-neighbor-witnesses}
\end{equation}
The magnitude bounds \(B,H\) are nonnegative.  The residual estimates
allow \(s=0\); the resource inequalities below that divide by \(s\) assume
\(s>0\).

\begin{theorem}[Single-flip convex mixture]
\label{thm:frontier-single-flip}
The average \(\bar x=N^{-1}\sum_i x^{(i)}\) is nonnegative and satisfies
\begin{equation}
 \boxed{\norm{A_{\varsigma}\bar x-b}_2
 \leq \frac{2BH\sqrt{sM_{\rm dep}}}{N}.}
 \label{eq:frontier-single-flip}
\end{equation}
If a public objective \(c\) has the same optimum \(v_*\) on every instance
and the witnesses in \eqref{eq:frontier-neighbor-witnesses} are optimal, then
\(c^\top\bar x=v_*\): the mixture is objective-exact, although it is
generally infeasible for the base instance.
\end{theorem}

\begin{proof}
Feasibility for the neighboring instances gives
\[
 N(A_{\varsigma}\bar x-b)
 =\sum_{i=1}^N
   (A_{\varsigma}-A_{\varsigma^{(i)}})x^{(i)}.
\]
Let \(D_r\) be the dependent positions in row \(r\), and set
\(m_r=|D_r|\).  In the \(r\)-th coordinate, position \((r,j)\) contributes
only to the summand \(i=\ell(r,j)\), so this coordinate is a sum of at most
\(m_r\) terms, each of magnitude at most \(2BH\).
Since \(m_r\leq s\),
\[
 \left|N(A_{\varsigma}\bar x-b)_r\right|^2
 \leq4B^2H^2m_r^2
 \leq4sB^2H^2m_r.
\]
Summing over rows and using \(\sum_rm_r=M_{\rm dep}\) proves
\eqref{eq:frontier-single-flip}.  Convexity proves nonnegativity, and
linearity proves the objective statement.
\end{proof}

The theorem also covers a right-hand side whose entries each depend on one
bit: apply it to \([A_{\varsigma},-b_{\varsigma}](x,1)=0\).  This uses row
sparsity \(s+1\) and height \(\max\{H,1\}\), replaces \(B\) by a bound on
every entry of \([A_{\varsigma},-b_{\varsigma}]\), and counts the dependent
RHS positions.  Inequalities are covered after slack conversion only when
the selected slack coordinates satisfy the same height bound.

The output hypothesis cannot be omitted.  For an affine public score
\(\phi(x)=a^\top x-a_0\), suppose every selected witness for input \(\xi\)
obeys \(p(\xi)\phi(x^\xi)\geq\gamma>0\).  All single-flip neighbors have
parity \(-p(\varsigma)\), so
\(p(\varsigma)\phi(\bar x)\leq-\gamma\).  Assume \(b\neq0\), as in
\cref{def:global-residual}, and call an output requirement
\emph{sound at threshold \(\eta\geq0\)} if it rejects every nonnegative,
objective-exact point with the wrong output and relative residual at most
\(\eta\).  A sound output requirement needs
\begin{equation}
 M_{\rm dep}B^2H^2
 >\frac{\eta^2\norm b_2^2}{4s}N^2.
 \label{eq:frontier-mixture-product}
\end{equation}
In particular, for fixed \(s,B,\eta,\norm b_2\), linear incidence and
bounded height are incompatible with soundness at a fixed threshold
\(\eta\) (a fixed residual tube).  Here and in later asymptotic statements
about a fixed tube, \(\eta>0\).

The same conclusion holds for an output requirement on the normalized primal
state \(x/\norm x_2\) under a concrete mixture-stability condition.  Let
\(\mathcal S\) be \(L\geq1\) pairwise disjoint public output pairs
\((u_\ell,v_\ell)\), and suppose
\begin{equation}
 p(\xi)(u_\ell^\xi-v_\ell^\xi)\geq a>0
 \quad(\ell\in\mathcal S)
 \label{eq:frontier-pair-margin}
\end{equation}
for every selected witness.  Averaging gives
\(-p(\varsigma)(\bar u_\ell-\bar v_\ell)\geq a\).  Nonnegativity gives
\(\bar u_\ell+\bar v_\ell\geq a\), hence
\[
 -p(\varsigma)(\bar u_\ell^2-\bar v_\ell^2)\geq a^2.
\]
For
\(Q_{\mathcal S}=\sum_{\ell\in\mathcal S}
 (\ket{u_\ell}\!\bra{u_\ell}-\ket{v_\ell}\!\bra{v_\ell})\),
\begin{equation}
 -p(\varsigma)
 \left\langle\frac{\bar x}{\norm{\bar x}_2}\middle|
 Q_{\mathcal S}\middle|\frac{\bar x}{\norm{\bar x}_2}\right\rangle
 \geq\frac{La^2}{nH^2}.
 \label{eq:frontier-pair-decoder}
\end{equation}
Measure a coordinate, output its pair sign on the selected pairs, and
output an independent fair sign otherwise.  The wrong-parity bias of this
binary measurement is half the signed expectation in
\eqref{eq:frontier-pair-decoder}.
Thus \(La^2=\Omega(nH^2)\) gives a constant wrong-parity bias.
Without an affine score, a common output template, the pair margin
\eqref{eq:frontier-pair-margin}, or another explicit mixture-stability
property, \cref{thm:frontier-single-flip} makes no claim about a
normalized quadratic decoder.

Two refinements show the roles of sensitivity and uneven incidence.  For a
Boolean function \(f\), fix a base input whose sensitive
set \(S_f\) has size \(k>0\), and average only over its sensitive neighbors.
If \(M_S\) counts the positions labelled by bits in \(S_f\), the same proof
gives
\begin{equation}
 \norm{A_{\varsigma}\bar x_S-b}_2
 \leq\frac{2BH\sqrt{sM_S}}{k}.
 \label{eq:frontier-sensitivity}
\end{equation}
Under the common-objective and mixture-stable output hypotheses, soundness
at threshold \(\eta\) therefore requires
\[
 M_SB^2H^2>
 \frac{\eta^2\norm b_2^2k^2}{4s}.
\]
More sharply, let \(M_i\) be the number of positions labelled by sensitive
bit \(i\).  For weights \(w_i\geq0\) with \(\sum_iw_i=1\), set
\(x_w=\sum_iw_ix^{(i)}\).  Cauchy--Schwarz in each row gives
\begin{equation}
 \norm{A_{\varsigma}x_w-b}_2^2
 \leq4sB^2H^2\sum_iM_iw_i^2.
 \label{eq:frontier-weighted-mixture}
\end{equation}
When every \(M_i>0\), the right side is minimized by
\(w_i=M_i^{-1}/\sum_hM_h^{-1}\), which gives
\[
 \norm{A_{\varsigma}x_w-b}_2
 \leq2BH\sqrt{\frac{s}{\sum_iM_i^{-1}}}.
\]
The corresponding soundness condition is
\[
 B^2H^2>
 \frac{\eta^2\norm b_2^2}{4s}\sum_iM_i^{-1}.
\]
If some \(M_i=0\), that neighbor's witness is exactly feasible for the base
matrix and already makes the output requirement unsound.  These mixtures preserve
every uniform affine or pair margin and the common objective.  The resource
that matters is therefore \(\sum_iM_i^{-1}\), not merely the total
incidence.

For a fixed relative tube and fixed \(s,\eta,\norm b_2\),
Table~\ref{tab:frontier-currencies} lists three constructions that attain
\eqref{eq:frontier-mixture-product} up to constants.
\begin{table}[t]
 \centering
 \small
 \caption{Constructions attaining the one-bit-local mixture bound
 \eqref{eq:frontier-mixture-product} up to constants, each by increasing one
 resource.}
 \label{tab:frontier-currencies}
 \begin{tabular}{@{}lccc@{}}
  \toprule
  construction & \(M_{\rm dep}\) & \(B\) & \(H\) \\
  \midrule
  bounded-height parallel replication
    & \(\Theta(N^2)\) & \(\Theta(1)\) & \(\Theta(1)\) \\
  direct row scaling
    & \(\Theta(N)\) & \(\Theta(\sqrt N)\) & \(\Theta(1)\) \\
  gain--plateau scaling
    & \(\Theta(N)\) & \(\Theta(1)\) & \(\Theta(\sqrt N)\) \\
  \bottomrule
 \end{tabular}
\end{table}
Every row has \(M_{\rm dep}B^2H^2=\Theta(N^2)\).  The parallel and
gain--plateau families of
\cref{thm:lp-parallel-path,thm:lp-gain-plateau} attain the
first and third rows.  Scaling the propagation rows by \(\sqrt N\)
attains the second without changing the feasible set; the cost appears in
the matrix normalization.  This sharpness concerns only these scales, not
the conditioning or QIPM cost of the formulations.

This is not an extension-complexity lower bound.  Carr--Konjevod give an
\(O(N)\)-size dynamic-programming extended formulation for the parity
polytope, and Ermel--Walter develop further flow formulations
\cite{carr2005-polyhedral-combinatorics,ermel2018-parity-polytopes-and-binarization}.
Those formulations do not guarantee the correct parity output at every point
of a fixed equality-residual tube.  Sparse LP decoding also has classical
fractional pseudocodewords and graph-cover descriptions
\cite{feldman2005-using-linear-programming-to-decode,
koetter2005-characterizations-of-pseudo-codewords-of,
vontobel2005-graph-cover-decoding}; these are conceptually related but do
not give the quantitative bound \eqref{eq:frontier-single-flip}.  In
particular, expander LP decoding may use an input-dependent objective,
whereas exact objective preservation above requires one common objective
and optimum.

\subsection{Extensions: multi-bit coefficients, KKT systems, cones, and centrality}
\label{sec:frontier-mixture-supplements}

If coefficient position \((r,j)\) may depend on a set
\(I_{rj}\subseteq[N]\) of input bits, define
\[
 d=\max_r\left|\bigcup_jI_{rj}\right|,
 \qquad d_{\rm coeff}=\max_{r,j}|I_{rj}|\leq d,
 \qquad M=\sum_{(r,j)}|I_{rj}|,
\]
and keep the row sparsity \(s\), coefficient bound \(B\), and witness
height \(H\).

\begin{proposition}[Multi-bit coefficients]
\label{prop:frontier-multibit}
The single-flip average obeys
\begin{equation}
 \norm{A_{\varsigma}\bar x-b}_2
 \leq\frac{2BH\sqrt{d_{\rm coeff}sM}}{N}
 \leq\frac{2BH\sqrt{dsM}}{N}.
 \label{eq:frontier-multibit}
\end{equation}
For fixed \(d_{\rm coeff},s\), the product bound
\(MB^2H^2=\Omega(N^2)\) still holds under the objective and output
assumptions above.  The first bound recovers the one-bit constant when every
coefficient depends on at most one bit, however many distinct bits appear in
a row.
\end{proposition}

\begin{proof}
Write \(f_{rji}=(A_{\varsigma,rj}-A_{\varsigma^{(i)},rj})x_j^{(i)}\).
It vanishes unless \(i\in I_{rj}\), and \(|f_{rji}|\leq2BH\).
Applying Cauchy--Schwarz first over the at most \(s\) dependent positions in
a row, and then over each coefficient's input set, gives
\[
 \left(\sum_j\sum_{i\in I_{rj}}f_{rji}\right)^2
 \leq s\sum_j |I_{rj}|\sum_{i\in I_{rj}}f_{rji}^2
 \leq4s d_{\rm coeff} B^2H^2\sum_j|I_{rj}|.
\]
Sum over rows and divide by \(N^2\).  The row-union estimate follows
from \(d_{\rm coeff}\leq d\).
\end{proof}

A simultaneous primal--dual version also holds.  Suppose \(A_\varsigma\)
has maximum row and column sparsities \(s_r,s_c\) and \(M\)
one-bit-dependent positions, and that the selected exact optimal triples
\((x^\xi,y^\xi,s^\xi)\) have all coordinates bounded in magnitude by \(H\)
and a common optimal value \(v\).  Split \(y=y^+-y^-\), stack the primal
feasibility and dual stationarity equations, and put
\[
 s_{\rm KKT}=\max\{s_r,2s_c+1\},
 \qquad B_{\rm KKT}=\max\{B,1\}.
\]
Each dependent position occurs once in the primal block and twice in the
split dual block.  \Cref{thm:frontier-single-flip} therefore gives
\begin{equation}
 \left\|
 \begin{bmatrix}
 A_\varsigma\bar x-b\\
 A_\varsigma^\top\bar y+\bar s-c
 \end{bmatrix}\right\|_2
 \leq\frac{2B_{\rm KKT}H\sqrt{3s_{\rm KKT}M}}{N}.
 \label{eq:frontier-kkt-stack}
\end{equation}
By linearity \(c^\top\bar x=b^\top\bar y=v\), so the algebraic duality gap
is exactly zero.  This covers output requirements combining feasibility, stationarity,
and gap requirements; the mixture does not preserve coordinatewise
complementarity.

The mixture proof uses only convexity of the domain.  Let \(\mathcal K\) be
a fixed convex set or cone, and suppose that for every selected input
\(\xi\in\{\varsigma,\varsigma^{(1)},\ldots,\varsigma^{(N)}\}\) there is a
witness \(x^\xi\in\mathcal K\) with
\(A_\xi x^\xi=b\), \(\norm{x^\xi}_\infty\leq H\) in one fixed basis, and
\(c^\top x^\xi=v\) for one public objective and common value.  If, in
addition, a fixed parity-neutral linear map satisfies
\(Lx^\xi=p(\xi)y\) for a fixed nonzero \(y\), then the half-base mixture
\[
 \widehat x=\frac12x^\varsigma+
              \frac1{2N}\sum_i x^{\varsigma^{(i)}}
\]
belongs to \(\mathcal K\), has the common objective, satisfies
\(L\widehat x=0\), and has half the residual in
\eqref{eq:frontier-multibit}.  For a sparse SDP this statement applies after
fixing one \(\operatorname{svec}\) identification, with all sparsity,
locality, coefficient, and coordinate bounds measured in that basis.  By
itself it gives no conclusion about normalized states, rank, determinant
neighborhoods, or complementarity.

A mixture of central points stays near centrality only under an explicit
stability hypothesis.  Keep the preceding row, column, one-bit-locality, and
height bounds, and suppose the selected neighboring triples are exact
central points for a common \(\mu>0\), so \(x_j^{(i)}s_j^{(i)}=\mu\).  Let
\(\bar x,\bar s\) be their averages.  For a fixed coordinate \(j\), write
\(t_i=x_j^{(i)}\).  Symmetrizing over pairs \(i,k\) gives the exact
identity
\begin{equation}
 \boxed{
 \frac{\bar x_j\bar s_j}{\mu}-1
 =\frac1{2N^2}\sum_{i,k=1}^N
   \frac{(t_i-t_k)^2}{t_it_k}.}
 \label{eq:frontier-multiplicative-variance}
\end{equation}
In particular, the excess is nonnegative and vanishes exactly when the
neighboring \(j\)-th primal coordinates agree.  If, for every coordinate
\(j\),
\(\max_i x_j^{(i)}/\min_i x_j^{(i)}\leq R\), where \(R\geq1\), then
\begin{equation}
 \mu\leq\bar x_j\bar s_j\leq K(R)\mu,
 \qquad K(R)=\frac{(R+1)^2}{4R}
 =1+\frac{(R-1)^2}{4R}.
 \label{eq:frontier-kantorovich}
\end{equation}
Indeed, for \(\alpha\leq t_i\leq\beta\),
\(t_i+\alpha\beta/t_i\leq\alpha+\beta\).  Averaging and maximizing the
product under this linear constraint gives
\((N^{-1}\sum_it_i)(N^{-1}\sum_it_i^{-1})
\leq(\alpha+\beta)^2/(4\alpha\beta)\), the scalar Kantorovich inequality
\cite{kantorovich1948-functional-analysis-applied-mathematics}.

\begin{proposition}[Central-mixture soundness dichotomy]
\label{prop:frontier-central-dichotomy}
Assume separately that the output is wrong for the base input uniformly
over mixtures, in one of two forms.  Write
\(Z^{(i)}=(x^{(i)},y^{(i)},s^{(i)})\).  Either a fixed affine
decoder \(\Phi\) satisfies
\(p(\varsigma)\Phi(Z^{(i)})\leq-\gamma<0\) for every \(i\), or the
primal components have the public pair decoder of
\eqref{eq:frontier-pair-decoder}, with
\(-p(\varsigma)(u_\ell^{(i)}-v_\ell^{(i)})\geq a\) for every \(i,\ell\)
and \(La^2=\Omega(nH^2)\).  In the first case the decoded object is the
stacked triple \(Z_w\); in the second it is the normalized primal state
\(\ket{x_w}=x_w/\norm{x_w}_2\).

Let \(w\) be any probability vector and set
\(Z_w=(x_w,y_w,s_w)=\sum_iw_i(x^{(i)},y^{(i)},s^{(i)})\).  With \(M_i\) the
one-bit incidence of bit \(i\), and with \(t_{ij}=x_j^{(i)}\), the stacked
residual and the centrality defect obey
\begin{align}
 \left\|
 \begin{bmatrix}A_\varsigma x_w-b\\
 A_\varsigma^\top y_w+s_w-c\end{bmatrix}\right\|_2^2
 &\leq12s_{\rm KKT}B_{\rm KKT}^2H^2\sum_iM_iw_i^2,
 \label{eq:frontier-weighted-kkt}\\
 \delta_j(w):=\frac{(x_w)_j(s_w)_j}{\mu}-1
 &=\frac12\sum_{i,k}w_iw_k
   \frac{(t_{ij}-t_{kj})^2}{t_{ij}t_{kj}}.
 \label{eq:frontier-weighted-centrality}
\end{align}
The weighted residual estimate admits the sharper constant
\begin{equation}
 \left\|
 \begin{bmatrix}A_\varsigma x_w-b\\
 A_\varsigma^\top y_w+s_w-c\end{bmatrix}\right\|_2^2
 \leq4(s_r+s_c)B^2H^2\sum_iM_iw_i^2.
 \label{eq:frontier-direct-kkt}
\end{equation}
Apply the weighted matrix-change estimate separately to the primal block
and to the transposed coefficient changes in stationarity.  The latter
allows signed multipliers and has the same incidence profile,
with column sparsity \(s_c\).  Adding the squared bounds proves
\eqref{eq:frontier-direct-kkt}; the earlier bound follows by increasing
the constant.  Thus the soundness inequalities below also hold with
\(12s_{\rm KKT}B_{\rm KKT}^2\) replaced by \(4(s_r+s_c)B^2\).
The algebraic objective difference, primal minus dual objective, remains
exactly \(n\mu\).  This follows from feasibility, stationarity, and
complementarity of each neighboring triple; the neighbors need not share
their individual primal or dual objective values.  An output requirement asking for
absolute stacked residual at most \(\varepsilon\) and an \(\ell_\infty\)
neighborhood of width \(\theta\) around the common \(\mu\) is therefore
sound only if, for every \(w\),
\begin{equation}
 12s_{\rm KKT}B_{\rm KKT}^2H^2\sum_iM_iw_i^2>\varepsilon^2
 \quad\text{or}\quad
 \max_j\delta_j(w)>\theta.
 \label{eq:frontier-central-dichotomy}
\end{equation}
Here \(\varepsilon,\theta\geq0\).  For the point-centered convention,
\(\mu_w=x_w^\top s_w/n\) with \(n>0\), replace the second term by the exact
width
\[
 \frac{\max_j|\delta_j-\bar\delta|}{1+\bar\delta},
 \qquad \bar\delta=\frac1n\sum_j\delta_j.
\]
For uniform weights, suppose for every coordinate \(j\) that
\[
 \frac{\max_i t_{ij}}{\min_i t_{ij}}\leq R.
\]
Then the width is at most \(K(R)-1\) in either convention, so soundness
requires the simpler dichotomy
\begin{equation}
 \frac{12s_{\rm KKT}B_{\rm KKT}^2H^2M}{N^2}>\varepsilon^2
 \quad\text{or}\quad
 \frac{(R-1)^2}{4R}>\theta.
 \label{eq:frontier-central-uniform}
\end{equation}
\end{proposition}

\begin{proof}
The stacked mixture estimate is the single-flip argument applied to primal
feasibility and stationarity.  Pairwise symmetrization gives the displayed
formula for \(\delta_j(w)\), and linearity preserves the algebraic objective
difference.  The affine or pair-margin hypothesis makes the decoded object wrong for the
base input.  If neither strict inequality in
\eqref{eq:frontier-central-dichotomy} held, this mixture would satisfy
every condition of the output requirement.  The coordinatewise Kantorovich
bound gives \eqref{eq:frontier-central-uniform} for uniform weights.
\end{proof}

The stability hypothesis is necessary.  Fix \(0<\mu<1\), let
\(b=\mu e,c=0\), and let bit \(j\) choose
\(a_j\in\{\mu,1\}\) in the diagonal LP
\(\operatorname{Diag}(a)x=b\).  Its exact \(\mu\)-central point is
\[
 x_j=\mu/a_j,\qquad s_j=a_j,\qquad y_j=-1.
\]
At the all-\(\mu\) base input, each single-flip neighbor changes one
coordinate from \((x,s)=(1,\mu)\) to \((\mu,1)\).  Thus every coordinate of
the average has product
\[
 \bar x_j\bar s_j
 =\frac{N-1+\mu}{N}\frac{1+(N-1)\mu}{N},
\]
which is \(\Theta(1/N)\), rather than \(\Theta(\mu)\), when
\(\mu\ll1/N\), although the primal and stationarity residuals vanish as
\(O(N^{-1/2})\).  Bounded coefficients and bounded central coordinates alone
therefore do not keep central mixtures near the common \(\mu\).  This
particular uniform mixture has the same complementarity
product in every coordinate, so its point-centered width is exactly zero.
Its algebraic objective difference is still \(N\mu\), whereas its
complementarity sum is \(N\bar x_j\bar s_j\); the two need not agree at
an infeasible mixture.

A bounded-coordinate counterexample also exists for the point-centered
convention.  Append \(N\) public diagonal coordinates with
coefficient \(\mu\) and central coordinates \((x,s,y)=(1,\mu,-1)\).
The moving coordinates share the product
\(q=(N-1+\mu)(1+(N-1)\mu)/N^2\), while the public coordinates have
product \(\mu\).  The point-centered mean is \((q+\mu)/2\), and
every coordinate has relative deviation \((q-\mu)/(q+\mu)\).
For \(N\geq8\) and \(\mu=N^{-2}\), the exact values satisfy
\[
 \left\|\begin{bmatrix}A\bar x-b\\A^\top\bar y+\bar s-c\end{bmatrix}\right\|_2^2
 =\frac{(1+\mu^2)(1-\mu)^2}{N}\leq\frac2N,
 \qquad \frac{q-\mu}{q+\mu}\geq\frac13.
\]
The unpadded example has the same residual and
\(\max_j\delta_j\geq N/8\) for \(\mu=N^{-2}\), \(N\geq2\).
Both constructions have row and column sparsity one, coefficient and
coordinate magnitudes at most one, and exactly \(N\) dependent positions.
Thus these bounds and a vanishing residual do not imply a point-centered
width below any fixed \(\theta<1/3\).

\paragraph{Formal verification.}
The Lean development in \path{formal/QipmFormal/Mixture/} verifies the
residual and centrality results of these two subsections from the actual
finite matrices and neighboring witnesses.  It covers weighted and
sensitive-subset mixtures, the stronger coefficientwise multibit bound,
RHS and slack augmentation, the split KKT construction, the direct KKT
bound, harmonic incidence weights, affine and normalized pair decoders,
the scalar variance and Kantorovich identities, both neighborhood
conventions, and the explicit counterexamples.  The soundness hypotheses
are stated as conditions on candidate outputs; the proofs derive the
residual, objective, and decoder properties of the mixture.  The development
uses Euclidean sums of squares and relies only on Lean's standard axioms
\texttt{propext}, \texttt{Classical.choice}, and \texttt{Quot.sound}.
The claim map and reproduction command are in \path{formal/MIXTURE.md}.
The sharpness constructions in Table~\ref{tab:frontier-currencies},
the electrical bounds below, and the quantum query reductions are not
covered by this verification.  The noncommutative extension has a separate
Lean development and claim map in \path{formal/SDP_MIXTURE.md}; see
\cref{sec:structural-sdp-mixtures}.

\subsection{Electrical-resistance bounds}
\label{sec:frontier-electrical}

The graph argument constructs, within the base instance, a point that
erases the outputs and a point that flips them, and shows the geometry behind the mixture bound: disjoint input
cuts force a large effective resistance.  Let
\(G=(V,E)\) be a connected oriented multigraph with root \(r\), nonempty
output set \(S\), public heights \(h_v>0\) with \(h_r=1\), and propagation
rows
\begin{equation}
 \lambda_e(d_v-g_ea_e(\varsigma)d_u)=0,
 \qquad g_e=h_v/h_u,
 \label{eq:frontier-signed-edge}
\end{equation}
where \(e=(u,v)\), \(\lambda_e\neq0\) is a public scale, \(a_e\) is either a
fixed sign or a fixed sign times one input bit, and the rows are
cycle-consistent.  Thus there are signs \(\tau_v\) with
\(\tau_r=1\) and \(\tau_v=a_e\tau_u\).  Assume separately that the
designated outputs carry the parity:
\begin{equation}
 \tau_s(\varsigma)=\epsilon_s p(\varsigma)
 \quad(s\in S),\qquad \epsilon_s\in\{\pm1\}\text{ public}.
 \label{eq:frontier-cut-output-parity}
\end{equation}
Embed each \(d_v\) using the capped pair of
\cref{sec:lp-foundation}: write \(d_v=u_v-v_v\) and
\(q_v=u_v+v_v\), fix a reference variable \(h_v^{\rm var}=h_v\), impose
\(q_v+t_v-2h_v^{\rm var}=0\), and use the local objective
\(2u_v+2v_v+h_v^{\rm var}+t_v\).  Put
\[
 L=\max_{e=(u,v)}|\lambda_e|h_v\leq BH,
 \qquad M=|E|,
\]
where the inequality assumes \(0<|\lambda_e|\leq B\) and \(h_v\leq H\).

\begin{theorem}[Cut--resistance bound]
\label{thm:frontier-cut-resistance}
There is a nonnegative, objective-exact point that satisfies every root,
reference, and cap row exactly, erases all output differences, and obeys
\begin{equation}
 \norm{Ax-b}_2^2\leq\frac{ML^2}{N^2}
 \leq\frac{MB^2H^2}{N^2}.
 \label{eq:frontier-cut-residual}
\end{equation}
There is also a point with every output sign flipped, whose squared residual
is at most four times this bound.  If the two unit anchors are the only
nonzero RHS entries, so \(\norm b_2=\sqrt2\), any relative-residual output requirement
with threshold \(\eta\) that rules out the erased point requires
\begin{equation}
 MB^2H^2>2\eta^2N^2.
 \label{eq:frontier-cut-product}
\end{equation}
If each edge label contains at most \(k\) raw signs, the residual bound
\eqref{eq:frontier-cut-residual} for the erased point acquires a factor
\(k^2\), and soundness requires \(Mk^2B^2H^2=\Omega(N^2)\).
\end{theorem}

\begin{proof}
For bit \(i\), let \(E_i\) be the edges whose label changes when bit
\(i\) flips.  By cycle consistency every cycle meets \(E_i\) evenly, so signed-graph switching identifies \(E_i\) with a cut
separating the root from every output
\cite{harary1953-balance-signed-graph,
zaslavsky2013-matrices-signed-simple-graphs}.  One-bit locality makes the
\(N\) cuts edge-disjoint.

Give edge \(e=(u,v)\) conductance
\(c_e=(\lambda_eh_v)^2\leq L^2\), and merge all outputs into one node.  For any
unit root-to-output flow \(f\), one unit crosses every cut.  Weighted
Cauchy--Schwarz gives
\[
 \sum_{e\in E_i}\frac{f_e^2}{c_e}\geq\frac1{C_i},
 \qquad C_i=\sum_{e\in E_i}c_e.
\]
Because the cuts are disjoint, Thomson's principle and a second application
of Cauchy--Schwarz give the weighted Nash--Williams bound
\begin{equation}
 \mathcal R_{r,S}\geq\sum_i\frac1{C_i}
 \geq\frac{N^2}{\sum_iC_i}
 \geq\frac{N^2}{ML^2}.
 \label{eq:frontier-nash-williams}
\end{equation}
We include this derivation to show where the oracle's disjoint input cuts
enter; the resistance inequality itself is classical
\cite{nashwilliams1959-random-walk-electric-currents}.

Let \(w\in[0,1]^V\) be the harmonic voltage with \(w_r=1\) and \(w_s=0\)
on every output.  By Dirichlet's principle and
\eqref{eq:frontier-nash-williams}, its energy is at most \(ML^2/N^2\).  Set
\[
 d_v=\tau_vh_vw_v,\quad q_v=h_v^{\rm var}=t_v=h_v,
 \quad u_v=(h_v+d_v)/2,\quad v_v=(h_v-d_v)/2.
\]
The maximum principle gives nonnegativity.  Since every $\lambda_e$ is
nonzero and the graph is connected, every feasible point has $|d_v|=h_v$,
hence $q_v=u_v+v_v\ge|d_v|=h_v$; the local objective reduces to $q_v+3h_v$,
so this point, which has $q_v=h_v$, is objective-exact.  All nonpropagation
rows hold exactly, while edge \(e\) has residual
\(\lambda_e\tau_vh_v(w_v-w_u)\), whose squared sum is the Dirichlet
energy.  Every output has \(d_s=0\).
Replacing the output boundary value zero by \(-1\) doubles every voltage
difference and proves the statement about the sign-flipped point.

For \(k\)-local labels, an edge lies in at most \(k\) of the bit cuts.  Hence
\(\mathcal R_{r,S}\geq k^{-1}\sum_iC_i^{-1}\), while
\(\sum_iC_i\leq kML^2\).  This gives
\(\mathcal R_{r,S}\geq N^2/(k^2ML^2)\) and the stated factor.
\end{proof}

By \cref{thm:frontier-cut-resistance}, the gain--plateau family is
optimal among cycle-consistent signed graphs: linear size and bounded row
scales force \(H=\Omega(\sqrt N)\).  The other two rows of
Table~\ref{tab:frontier-currencies} attain the remaining extremes.
Adaptivity does not rescue the pointwise output requirement, because the
harmonic point exists at every requested stage.  However, the theorem gives
no unconditional lower bound on fresh queries per iteration:
after all \(N\) fixed signs are cached, later raw-query cost may be zero
under \cref{def:amortized-accounting}.

The resistance argument extends to rows with several wires.
Give wire \(v\) a public height \(h_v\in(0,H]\) and intended value
\(h_v\chi_{J_v}(\varsigma)\), where
\(\chi_J=\prod_{i\in J}\varsigma_i\).  Let each of \(m\) homogeneous rows
have arity at most \(q\):
\begin{equation}
 \sum_{v\in V_a}\alpha_{av}\chi_{I_{av}}(\varsigma)d_v=0,
 \qquad |\alpha_{av}|\leq A,\qquad |I_{av}|\leq k.
 \label{eq:frontier-walsh-row}
\end{equation}
Roots have \(J_v=\varnothing\), outputs have \(J_v=[N]\), and every row
holds on the intended wires for every input.  We may assume that the
comparison graph in the proof connects the roots to the outputs; otherwise the output character is unconstrained and can be erased with zero
propagation residual.
For the objective-exact conclusion, assume also that the intended
full-height assignment is optimal for the capped LP.  It suffices that
every exactly feasible assignment has \(|d_v|\geq h_v\) on every wire,
since nonnegativity then forces \(q_v\geq h_v\).

\begin{theorem}[Bounded-arity Walsh-row bound]
\label{thm:frontier-walsh}
The capped-pair embedding has a nonnegative, objective-exact point that
satisfies the anchors, references, and caps, erases every output, and has
\begin{equation}
 \norm r_2^2\leq
 \frac{4k^2q^3mA^2H^2}{N^2}.
 \label{eq:frontier-walsh-residual}
\end{equation}
Consequently, for fixed \(k,q\), soundness at a constant absolute residual
threshold requires \(mA^2H^2=\Omega(N^2)\).
\end{theorem}

\begin{proof}
Substituting the intended wires into row \(a\) gives
\(\sum_v\beta_{av}\chi_{K_{av}}=0\), where
\(\beta_{av}=\alpha_{av}h_v\) and
\(K_{av}=I_{av}\mathbin\triangle J_v\).  Since distinct Walsh characters
are linearly independent, \(\sum_{v\in G_{a,K}}\beta_{av}=0\) within each
equal-character group.  Choose one pivot per group and join it to the other
group members by comparison edges of conductance
\(q^2\beta_{av}^2\).  There are at most
\(qm\) edges, each with conductance at most \(q^2A^2H^2\).  For arbitrary
interpolation values \(w_v\), a character group with pivot \(p\) has
residual
\[
 \sum_{v\in G}\beta_{av}w_v
 =\sum_{v\in G\setminus\{p\}}\beta_{av}(w_v-w_p).
\]
Applying Cauchy--Schwarz first over the at most \(q\) terms of a group and
then over the at most \(q\) groups of a row gives
\[
 \norm r_2^2\leq
 \sum_ec_e(w_{h(e)}-w_{t(e)})^2.
\]

The endpoints of a comparison edge satisfy
\(J_u\triangle J_v=I_{au}\triangle I_{av}\), so the edge crosses at most
\(2k\) of the vertex cuts \(\{v:i\in J_v\}\).  The bounded-overlap form of
the Nash--Williams calculation proved above gives
\[
 \mathcal R_{\rm eff}
 \geq\frac{N^2}{4k^2q^3mA^2H^2}.
\]
Take the harmonic interpolation equal to one at the roots and zero at the
outputs.  Its energy gives \eqref{eq:frontier-walsh-residual}; the maximum
principle and the capped-pair construction from the preceding proof give
nonnegativity and exact auxiliary rows.  The erased point has
\(q_v=h_v\) on every wire, hence the local objective of the intended
assignment, which is optimal by hypothesis.
\end{proof}

Formulations built from exact local XOR hulls have three further concrete
defects.  If a set \(U\) of unfixed variables meets every parity check in zero or at
least two positions, then the point that sets \(x_j=1/2\) on \(U\) and keeps
a satisfying Boolean codeword elsewhere lies in every local parity polytope.
Indeed, for a check meeting \(U\) in \(r\geq2\) variables, this point is the
uniform average of the \(2^{r-1}\) assignments with the required residual
parity.  This stopping-set pseudocodeword can make an output exactly
unbiased.  If the checks instead peel along a dependency chain of length
\(D\), linearly interpolating the gauged wire error from zero to one
violates one opposing facet at each gate by \(1/D\), for total Euclidean
violation \(1/\sqrt D\).  Finally, every satisfying truth-table assignment
is a vertex of its exact gate hull, so some standard-form slacks vanish; a
Boolean circuit forced by its inputs is not strictly primal feasible.
Thickening the hulls restores positive slacks only by admitting fractional
drift.  These facts explain the LP-decoding analogy but are not a universal
lower bound for all inequality extended formulations.

\subsection{Multiplier conservation}
\label{sec:frontier-multiplier}

Residual robustness constrains primal witnesses.  Complete central and
Newton outputs face a second constraint: conservation of the dual
multipliers.  In a signed-copy graph \(G=(V\cup T,E)\), vertices in \(V\) carry variables with
\(d_v=x_v^+-x_v^-\) and \(q_v=x_v^++x_v^-\), and vertices in \(T\) are fixed
terminals used for anchor rows.  A variable--variable edge \(e=(u,v)\)
represents \(d_v-a_ed_u=0\).  A terminal edge is oriented from its terminal
to its variable endpoint, with the fixed terminal term moved to the RHS.
Let \(B_a\) and \(B\) be the signed and the ordinary oriented incidence
matrices, with all terminal columns omitted.  For
\(S\subseteq V\), the boundary \(\partial S\) consists of the edges to
\(V\setminus S\) and all terminal edges incident to \(S\).  Each normalized
edge equation has a multiplier \(y_e\); an actual LP row may be multiplied
by a public nonzero scale \(r_e\ne0\), in which case its scale-invariant row
force is \(f_e=r_ey_e\).
Assume exact primal--dual centrality, \(d_v\ne0\), a pair-symmetric objective,
and pair-symmetric columns for all equality constraints outside the graph.

\begin{theorem}[Gauge flow and isoperimetric law]
\label{thm:frontier-gauge-flow}
Let \(\tau_v=\operatorname{sgn}(d_v)\).  After the edge gauge
\(z_e=\tau_{h(e)}y_e\), signed stationarity becomes the positive-divergence
law
\begin{equation}
 B^\top z=g,
 \qquad g_v=\frac{2\mu|d_v|}{q_v^2-d_v^2}>0.
 \label{eq:frontier-gauge-flow}
\end{equation}
For every vertex set \(S\) with nonempty boundary,
\begin{equation}
 \sum_{e\in\partial S}|y_e|\geq G_S,
 \quad
 \norm{y_{\partial S}}_\infty\geq\frac{G_S}{|\partial S|},
 \quad
 \norm{y_{\partial S}}_2\geq\frac{G_S}{\sqrt{|\partial S|}},
 \quad G_S=\sum_{v\in S}g_v.
 \label{eq:frontier-cut-flow}
\end{equation}
If \(S\ne\varnothing\) but \(\partial S=\varnothing\), no such positive
central point exists.
If \(|d_v|\geq a>0\) and \(q_v\leq Q\) on \(S\), then
\begin{equation}
 \norm{y_{\partial S}}_\infty
 \geq\frac{2\mu a}{Q^2-a^2}\frac{|S|}{|\partial S|},
 \qquad
 \norm{y_{\partial S}}_2
 \geq\frac{2\mu a}{Q^2-a^2}\frac{|S|}{\sqrt{|\partial S|}}.
 \label{eq:frontier-isoperimetric}
\end{equation}
For a row written as \(r_e(d_v-a_ed_u)=0\), all statements hold for the
row force \(f_e=r_ey_e\); hence, if \(|r_e|\leq B_{\rm row}\), the lower
bounds apply to \(B_{\rm row}\norm y_\infty\), not to the multiplier
alone, which depends on the row normalization.
\end{theorem}

\begin{proof}
Subtract the two stationarity equations of each pair:
\[
 2(B_a^\top y)_v+s_v^+-s_v^-=0.
\]
Since \(x_v^\pm=(q_v\pm d_v)/2\) and
\(s_v^\pm=\mu/x_v^\pm\), this becomes
\((B_a^\top y)_v=2\mu d_v/(q_v^2-d_v^2)\).  Feasibility implies
\(a_e=\tau_u\tau_v\) on variable--variable edges; terminal rows have the
same single variable coefficient in \(B_a\) and \(B\).  Switching the signed
incidence matrix therefore gives
\eqref{eq:frontier-gauge-flow}.  Summing over \(S\) cancels the internal
edges; the triangle inequality and Cauchy--Schwarz give
\eqref{eq:frontier-cut-flow}.  Positivity gives \(q_v>|d_v|\), and the
margin hypotheses bound every \(g_v\) below by
\(2\mu a/(Q^2-a^2)\).  A diagonal row scaling enters stationarity only
through \(r_ey_e\), which proves the final statement.
\end{proof}

Pair symmetry is not needed to compare two inputs if the pair-antisymmetric
source \(h_v\) below is the same vector, in the original pair coordinates,
for both.  Let two exact central points, denoted \(+\) and \(-\), have
row forces \(f_e^\pm=r_ey_e^\pm\), and suppose their subtracted stationarity
equations are
\[
 (B_{a^\pm}^\top f^\pm)_v
 =h_v+\frac{2\mu_\pm d_v^\pm}
              {(q_v^\pm)^2-(d_v^\pm)^2}.
\]
If \(d_v^+>0>d_v^-\) on \(S\), then switching the two inputs, adding their
equations, and summing over \(S\) cancels both the common source and every
internal edge.  Hence
\begin{equation}
 \sum_{e\in\partial S}(|f_e^+|+|f_e^-|)
 \geq\sum_{v\in S}(g_v^++g_v^-),
 \qquad
 g_v^\pm=\frac{2\mu_\pm|d_v^\pm|}
 {(q_v^\pm)^2-(d_v^\pm)^2}.
 \label{eq:frontier-two-input}
\end{equation}
If \(\mu_\pm\geq\mu_0\), \(|d_v^\pm|\geq a\), \(q_v^\pm\leq Q\), and
\(|r_e|\leq B_{\rm row}\), this gives
\begin{equation}
 \max_{p\in\{+,-\}}\norm{y_{\partial S}^p}_\infty
 \geq\frac{2\mu_0a}{Q^2-a^2}
       \frac{|S|}{B_{\rm row}|\partial S|}.
 \label{eq:frontier-two-input-norm}
\end{equation}
Thus a public asymmetric objective can cancel the demand \(g_v\) for one
sign pattern of \(d\) on \(S\), but not for both when \(\partial S\) is
small.  Coefficients of another public asymmetric constraint do not suffice
unless their contribution through the multipliers is also the same vector
\(h\) for both inputs.

The identity for the work \(b_a^\top y_a\) does not need a graph.  Let
\(Cd=b_a\) be an arbitrary pair-antisymmetric equality block, with columns
\(C,-C\), and let all other local data be pair-symmetric.

\begin{theorem}[Central and Newton work identities]
\label{thm:frontier-work}
At an exact central point,
\begin{equation}
 \boxed{b_a^\top y_a
 =\sum_v\frac{2\mu d_v^2}{q_v^2-d_v^2}\geq0.}
 \label{eq:frontier-central-work}
\end{equation}
The work is strictly positive exactly when \(d\ne0\).  At a pair-symmetric
Newton start with pair-symmetric dual and complementarity right-hand sides,
pair-symmetric other equality columns,
\(x_v^{0,+}=x_v^{0,-}=\xi_v\), and
\(s_v^{0,+}=s_v^{0,-}=\zeta_v\), the antisymmetric block satisfies
\begin{equation}
 C_\varsigma\Delta d=r_a,
 \qquad C_\varsigma^\top\alpha=D\Delta d,
 \qquad D=\operatorname{Diag}\!\left(\frac{\zeta_v}{2\xi_v}\right),
 \quad
 \boxed{r_a^\top\alpha=\Delta d^\top D\Delta d.}
 \label{eq:frontier-newton-work}
\end{equation}
Both pairings are invariant under
\((C,b_a,y_a)\mapsto(UC,Ub_a,U^{-\top}y_a)\), and likewise for
\((C_\varsigma,r_a,\alpha)\).  If \(D\succeq\rho I\),
\(|\Delta d_v|\geq a\) on \(L\) coordinates, and \(\norm{r_a}_1\leq R\),
then \(\norm\alpha_\infty\geq\rho a^2L/R\).
Likewise, if \(|d_v|\geq a\) and \(q_v\leq Q\) on \(L\) coordinates,
then
\begin{equation}
 \norm{y_a}_\infty\geq
 \frac{2\mu a^2L}{(Q^2-a^2)\norm{b_a}_1},
 \qquad
 \norm{y_a}_2\geq
 \frac{2\mu a^2L}{(Q^2-a^2)\norm{b_a}_2}.
 \label{eq:frontier-central-work-norms}
\end{equation}
A zero denominator with nonzero work means that no such point exists.
\end{theorem}

\begin{proof}
Pair subtraction gives
\(C^\top y_a=2\mu d/(q^2-d^2)\) coordinatewise.  Taking the inner product
with \(d\) and using \(Cd=b_a\) proves
\eqref{eq:frontier-central-work}.  At the Newton start, subtracting the dual
and complementarity equations and eliminating the slack difference gives
\(C_\varsigma^\top\alpha=D\Delta d\); primal antisymmetric feasibility gives
\(C_\varsigma\Delta d=r_a\).  Taking inner products proves
\eqref{eq:frontier-newton-work}, and H\"older's inequality gives the norm
bounds.  The contragredient transformations preserve both scalar pairings.
\end{proof}

For the bounded secant start of \cref{sec:lp-dual-homogeneity},
\(D=(2/9)I\), \(\Delta d_v=\tau_v\), and the antisymmetric RHS at the root
is a unit vector.  Equation~\eqref{eq:frontier-newton-work} therefore gives
the root multiplier \(2P/9\).  Scaling all dual data by \(\lambda\)
changes this to \(2\lambda P/9\); bounded full Newton directions require
\(\lambda=O(1/P)\), exactly the inverse dual scale used in
\cref{thm:lp-dual-homogeneous-full}.  \Cref{sec:lp-beyond-pairs}
states the corresponding unrestricted Newton virtual-work identity.  We use
the paired forms here because they expose the gadget resources;
\cref{sec:structural-tools} places them in the general
involution/symmetry framework.

The cut laws also hold directly for Newton directions.  Under a
pair-symmetric start, pair-symmetric Newton right-hand sides, and homogeneous
signed-copy primal equations, switching by
\(\operatorname{sgn}(\Delta d_v)\) gives
\begin{equation}
 B^\top z=g^{\rm N},
 \qquad g_v^{\rm N}=\frac{\zeta_v}{2\xi_v}|\Delta d_v|.
 \label{eq:frontier-newton-flow}
\end{equation}
The proof is the Newton half of \cref{thm:frontier-work} followed by
the cut summation in \cref{thm:frontier-gauge-flow}.  In particular,
if \(\zeta_v/\xi_v\geq r\) and \(|\Delta d_v|\geq a\) on \(S\), then
\[
 \norm{\Delta y_{\partial S}}_\infty
 \geq\frac{ra}{2}\frac{|S|}{|\partial S|}.
\]

Now assume that terminal values, labels, and terminal Newton residuals are
public and fixed, and fix a base input such that it and all its
single-flip neighbors have exact central points with nonzero \(d_v\).  Let
\(O\) be a fixed set of \(K\) output vertices on which
\(\operatorname{sgn}d_v\) flips under every single-bit flip, and define
\[
 S_j=\{v:\operatorname{sgn}d_v(\varsigma^{(j)})
             =-\operatorname{sgn}d_v(\varsigma)\}.
\]
Let \(M_j\) count the signed rows that change under flip \(j\), and put
\(M_{\rm dep}=\sum_jM_j\).  Feasibility on an internal edge gives
\(a_e=\operatorname{sgn}(d_u)\operatorname{sgn}(d_v)\), so its label changes
exactly when one endpoint is in \(S_j\).  Public terminal signs act as fixed
outside vertices, so the same boundary count includes terminal edges.
Hence
\(|\partial S_j|\leq M_j\), with equality when rows rather than expanded
coefficient positions are counted, and \(O\subseteq S_j\).  The base
multiplier must satisfy all \(N\) cut inequalities simultaneously.  Assume
explicitly that the base central point has
\(|d_v|\geq a=\Omega(1)\) and \(q_v\leq Q=O(1)\) on \(O\).  Applying this
central margin only on \(O\) and summing over \(j\) proves
\begin{equation}
 M_{\rm dep}B_{\rm row}\norm{y}_\infty
 \geq\frac{2\mu a}{Q^2-a^2}KN.
 \label{eq:frontier-central-incidence}
\end{equation}
For Newton directions, suppose instead that the base and every single-flip
system satisfy the hypotheses of \eqref{eq:frontier-newton-flow}, define
\(S_j\) from the signs of \(\Delta d\), and assume that every output
direction flips with \(|\Delta d_v|\geq a=\Omega(1)\) on \(O\).  The same
argument gives
\begin{equation}
 M_{\rm dep}B_{\rm row}\norm{\Delta y}_\infty
 \geq\rho aKN,
 \qquad \rho=\min_{v\in O}\frac{s_v^0}{2x_v^0}.
 \label{eq:frontier-newton-incidence}
\end{equation}
Thus \(M_{\rm dep}=\Theta(N)\), \(K=\Theta(N)\), bounded row scale, and
bounded multipliers, together with the explicit constant margin
\(a=\Omega(1)\) and central range \(Q=O(1)\), force
\(\mu=O(1/N)\) in the central system and \(\rho=O(1/N)\) in the Newton
system.  Equation
\eqref{eq:frontier-newton-incidence} is the general conservation law behind
the multiplier lower bound beyond signed pairs in \cref{sec:lp-beyond-pairs}.
Every signed-copy gadget must make the input incidence, row scale,
multiplier size, or inverse dual scale large.  Branching and redundant copy
rows only move the cost into the relevant boundary count.

Pair-asymmetric local data do not provide a cheap universal cancellation.
Suppose a prescribed direct source \(h_v(\varsigma)\) changes central
stationarity to
\[
 (B_a^\top f)_v=
 \frac{2\mu d_v}{q_v^2-d_v^2}+h_v(\varsigma),
 \qquad f_e=r_ey_e.
\]
On a fixed set \(O\) of \(K\) output vertices, assume
\(\operatorname{sgn}(d_v)=\kappa_vp(\varsigma)\),
\(|d_v|\geq a\), and \(q_v\leq Q\) for every input.  With the uniform
Fourier convention
\(\widehat h_v([N])=2^{-N}\sum_\varsigma p(\varsigma)h_v(\varsigma)\), put
\(\varepsilon_O=K^{-1}\sum_{v\in O}|\widehat h_v([N])|\).

\begin{theorem}[Fourier obstruction to direct cancellation]
\label{thm:frontier-fourier-source}
For some input \(\varsigma^*\),
\begin{equation}
 \sum_{e\in\partial O}|r_ey_e(\varsigma^*)|
 \geq K\left(\frac{2\mu a}{Q^2-a^2}-\varepsilon_O\right).
 \label{eq:frontier-fourier-source}
\end{equation}
In particular, bounded-bit juntas that depend on strictly fewer than \(N\)
input bits, and all sources of Fourier degree less than \(N\), have zero
full-parity coefficient and cannot cancel the demand for all inputs.  The
bound is tight: the direct source
\(h_v=-\kappa_vp(\varsigma)g_0\) cancels a constant demand \(g_0\) exactly.
\end{theorem}

\begin{proof}
After switching, sum the stationarity equations over \(O\).  The direct
contribution is
\(H(\varsigma)=\sum_{v\in O}\kappa_vp(\varsigma)h_v(\varsigma)\), whose
average over inputs is at least \(-K\varepsilon_O\).  Some \(\varsigma^*\) has
\(H(\varsigma^*)\) at least this average.  The positive central demand is at
least \(2\mu aK/(Q^2-a^2)\); internal edge forces cancel, and the triangle
inequality proves \eqref{eq:frontier-fourier-source}.
\end{proof}

There is an exact Newton analogue.  At a pair-symmetric start put
\(x_v^{0,\pm}=\xi_v\), \(s_v^{0,\pm}=\zeta_v\), and allow asymmetric dual
and complementarity right-hand sides.  Pair subtraction gives
\[
 (B_a^\top\Delta f)_v
 =\frac{\zeta_v}{2\xi_v}\Delta d_v+h_v^{\rm N}(\varsigma),
 \qquad
 h_v^{\rm N}=\frac{r_{d,v}^+-r_{d,v}^-}{2}
 -\frac{r_{c,v}^+-r_{c,v}^-}{2\xi_v}.
\]
If \(\operatorname{sgn}(\Delta d_v)=\kappa_vp(\varsigma)\),
\(|\Delta d_v|\geq a\), and \(\zeta_v/(2\xi_v)\geq\rho\) on \(O\), then
the same proof gives some input \(\varsigma^*\) with
\begin{equation}
 \sum_{e\in\partial O}|r_e\Delta y_e(\varsigma^*)|
 \geq K(\rho a-\varepsilon_O^{\rm N}),
 \qquad
 \varepsilon_O^{\rm N}=K^{-1}\sum_{v\in O}
 |\widehat h_v^{\rm N}([N])|.
 \label{eq:frontier-newton-fourier}
\end{equation}
This Fourier coefficient must be computed after division by the start
coordinates: input dependence in \(1/\xi_v\) can create full parity from
raw data that are otherwise local.  The word \emph{direct} also matters.
An unknown multiplier of another antisymmetric constraint may depend on all
input bits even when that row is local, so
\cref{thm:frontier-fourier-source} does not assume that such a
multiplier has low Fourier degree.

Switching and conservation of incidence flows are classical
\cite{harary1953-balance-signed-graph,zaslavsky2013-matrices-signed-simple-graphs};
the work identities are special cases of standard KKT virtual-work and
Schur-energy identities.  Electrical flows on interior-point central paths
are also established algorithmic tools \cite{madry2013-navigating-central-path}.
The only contribution claimed here is that the local input cuts, the
parity-sized output region, multiplier conservation, row scaling, and
incidence accounting must all be satisfied at once.  Effective resistance has a separate established role in
quantum graph algorithms
\cite{belovs2012-span-programs-st-connectivity,
jarret2018-quantum-connectivity,apers2022-quantum-speedup-for-graph-sparsification};
those works do not imply the LP residual or multiplier bounds of this
section.

\subsection{Relative-value estimation and long shortcut rows}
\label{sec:frontier-remnants}

The path family also gives a simple lower bound for value outputs.  Put
\(p_0=1\) and \(p_i=\prod_{k=1}^i\varsigma_k\).  On a single signed path,
minimize
\begin{equation}
 \frac1{N+1}\sum_{i=0}^N(u_i+v_i)
 +\frac14(u_N-v_N)
 \label{eq:frontier-relative-value-lp}
\end{equation}
subject to \(d_0=1\) and
\(d_i-\varsigma_id_{i-1}=0\).

\begin{theorem}[Relative-value gap]
\label{thm:frontier-relative-value}
In raw sparse coefficient access, estimating the optimum of
\eqref{eq:frontier-relative-value-lp} to target-relative error
\(\epsilon<1/4\) requires \(\Omega(N)\) queries.  The LP has a finite,
attained optimum, is strictly primal--dual feasible, and has a unique,
nondegenerate, strictly complementary optimal solution.  Its objective norm
is bounded by an absolute constant, and its primal log-barrier Hessian,
restricted to the feasible affine space, has condition number one along the
entire central path.
\end{theorem}

\begin{proof}
Exact feasibility fixes \(d_i=p_i\) and \(q_i=u_i+v_i\geq1\).  On this
affine space the objective is
\((N+1)^{-1}\sum_iq_i+p_N/4\), so the unique optimum has every \(q_i=1\)
and
\begin{equation}
 \operatorname{OPT}=1+\frac14p(\varsigma)
 \in\left\{\frac34,\frac54\right\}.
 \label{eq:frontier-relative-value-gap}
\end{equation}
The two target-relative error intervals are disjoint for every
\(\epsilon<1/4\); hence such value estimation requires \(\Omega(N)\) raw
queries by the parity lower bound.  Even an exactly feasible point with
objective at most \(3\operatorname{OPT}/2\) distinguishes the two cases by
a threshold between \(9/8\) and \(5/4\).  Each
sparse coefficient query is simulated by one raw sign query.  Also
\(\norm c_2^2=2/(N+1)+1/8=O(1)\).

The equality matrix is full row rank, and its public orthonormal null basis
consists of \((e_{u_i}+e_{v_i})/\sqrt2\).  Taking every \(q_i=2\) proves
strict primal feasibility.  Taking \(s_{u_i}=s_{v_i}=1/(N+1)\) gives a
positive slack vector for which \(c-s\) is the endpoint difference term; it is
orthogonal to the nullspace and hence lies in \(\operatorname{range}(A^\top)\),
proving strict dual feasibility.  At the optimum, one positive coordinate
from each pair gives a triangular basis, and every opposite coordinate has
positive reduced cost \(2/(N+1)\).  This proves nondegeneracy and strict
complementarity.

Finally, on the feasible slice the barrier objective is
\[
 \sum_{i=0}^N\left[\frac{q_i}{N+1}
 -\mu\log\frac{q_i^2-1}{4}\right]+\frac14p_N.
\]
Every central coordinate solves the same scalar equation
\((N+1)^{-1}=2\mu q/(q^2-1)\), so the Hessian in the displayed orthonormal
null basis is a positive scalar multiple of the identity.
\end{proof}

This is a value reduction with the normalization guarantees stated above,
not a query lower bound for generic LPs.

One tempting repair is a long two-coordinate shortcut row.  Suppose
\(0\leq i<j\leq N\) and a homogeneous row
\(a(\varsigma)d_i+b(\varsigma)d_j=0\) holds on every exact signed path,
with both \(a(\varsigma)\) and \(b(\varsigma)\) nonzero for every input
covered by the conclusion.  Substitution then gives
\begin{equation}
 -\frac{a(\varsigma)}{b(\varsigma)}
 =\frac{p_j}{p_i}=\prod_{k=i+1}^j\varsigma_k.
 \label{eq:frontier-shortcut-parity}
\end{equation}
Thus the row itself carries the parity of bits \(i+1,\ldots,j\),
invariant under common row rescaling, and generally costs
\(\Omega(j-i)\) raw queries to simulate.  Otherwise, replicated-path and
expander constructions are covered by the mixture, resistance, and
Walsh-row theorems above.

\subsection{Scope of the conclusions}
\label{sec:frontier-synthesis}

The parallel-path and gain--plateau constructions of \cref{sec:lp-hardness}
are optimal in their stated bounded-height and linear-size classes,
respectively.  The \(\Theta(N^{-1/2})\) residual threshold of
the all-unit family in \cref{thm:lp-unit-threshold} is the matching
threshold without amplification, and dual homogeneity attains the scaling
\(\rho=\Theta(1/N)\) that \eqref{eq:frontier-newton-incidence} forces.
None of these conclusions says that parity needs a quadratic-size LP, that
every bounded-range LP fails, or that an adaptive QIPM pays a fresh linear
query cost at every iteration.

A bounded-range construction with linear incidence must violate at least
one hypothesis used here.  The escape routes are:
\begin{enumerate}[label=(\roman*)]
 \item use an output promise that is nonlinear or otherwise not
       mixture-stable, and prove its robustness directly;
 \item require only central or Newton outputs, and allow the neighboring
       central coordinates to leave every fixed Kantorovich neighborhood
       \eqref{eq:frontier-kantorovich};
 \item use nonlocal input coefficients and charge their full raw-query
       cost;
 \item give up the common objective or common optimum, so neighbor
       mixtures need not be objective-exact; or
 \item impose integrality or another nonconvex promise, which leaves linear
       programming.
\end{enumerate}
The second route is not automatic:
\eqref{eq:frontier-central-dichotomy}--\eqref{eq:frontier-central-uniform}
show exactly how much central instability it requires.

Two problems remain open within the present scope.  First, can an arbitrary
sparse bounded-range equality LP with a nonlinear normalized-state decoder
be robust in a fixed global residual tube without nonlocal coefficient
access?  Second, can the bounded-scale full-KKT problem of
\cref{open:lp-bounded-full-kkt} be solved by a nonorthogonal
mixed-mode gadget when all input-dependent sources and recovery maps are
charged?  Together with the orthogonal-copy cut law in
\cref{sec:lp-beyond-pairs}, the theorems here rule out the natural
mixture-stable, signed-copy, constant-arity Walsh, permutation, rotation,
and flat-copy approaches; they do not settle either question.
